\begin{claim}{D.2.19}{Proved}{The next target group and the eight-skeleton extension}
Retain the notation of D.2.17--D.2.18. Then
\[
 \pi_7T=\pi_9(\Sigma^2K(\FF_2,2))\cong\ZZ/2,
 \qquad \rho\eta_6=0.
 \tag{D.2.19a}
\]
The generator has mod-two Hurewicz image $R=\{a,c\}$ in
$H_7T$. If $p:M(\ZZ/4,5)\to S^6$ is the top pinch,
$\zeta:S^7\to M(\ZZ/4,5)$ satisfies $p\zeta=\eta_6$,
and $\mu:M(\ZZ/4,5)\to T$ is the Moore extension used in
D.2.17, then
\[
 \mu\zeta=0.
 \tag{D.2.19b}
\]
The ordinary map on the seven-skeleton in D.2.18 can be chosen
to extend over a finite eight-skeleton $K_8$ of a CW model of
$A=Q\Sigma K(\FF_2,2)$, retaining all the homology values
in D.2.17. This is a finite-stage extension; no extension
over the full source or nullification conclusion is asserted.
\end{claim}
\begin{proof}
\textbf{The target degree is still stable.}
Let $F=\fib(T\to B)$ be the three-connected fiber of D.2.18.
It is a double loop space, since its defining map is the
double loop of $Y\to K(\FF_2,4)$. Write $F=\Omega Z$.
The adjunction triangle identity gives a left inverse
\[
 \Omega\epsilon_Z:\Omega\Sigma F\longrightarrow F
\]
to the suspension--loop unit. Thus that unit is injective
on every homotopy group. Freudenthal makes it surjective
on $\pi_7$, and $\Sigma F$ is four-connected. Hence
\[
 \pi_7F\cong\pi_8\Sigma F\cong\pi_7\Sigma^\infty F.
 \tag{D.2.19c}
\]
This removes the possible first unstable kernel, including
its torsion, without assuming that a proposed list of
Whitehead products exhausts it. The same retraction is
described by \cite[Theorem~1.8, p.~173]{I_James1955}.
We compute two-locally, as in D.2.18; the resulting finite
groups are also the integral groups.

\medskip\noindent
\textbf{Every degree-eight fiber class and its operations.}
Use the fiber classes
\[
 A_4=a^2,\quad z_5=e+ab,\quad u_5=\xi a,\quad
 d_6=b^2,\quad e_6=b^2+k,\quad f_6=p_3,
 \qquad k=\{a,b\}.
\]
Let $i,j,l\in H_8B$ be dual to $u^4,uv^2,vw$.
The full free double-loop basis of $H_8T$ is
\[
\begin{gathered}
 i,j,l,\ af,ax,bd,be,c^2,a^2c,ab^2,a^4,\\
 b\xi a,ak,\{a,d\},\{a,e\},\{b,c\}.
\end{gathered}
 \tag{D.2.19d}
\]
There is no further top-operation term in even degree eight.
The restricted Browder identities give
$\{a,\xi a\}=0$ and zero brackets on identical inputs;
these are the basis conventions of
\cite[Chapter~III, Theorems~1.2--1.4]{I_CohenLadaMay1976}.
The splitting $T\simeq B\times F$ subtracts dimensions
$3,1,2,3$ from the base and the mixed terms, leaving
$\dim H_8F=7$. A complete basis is
\[
\begin{split}
 W_8&=a^4,\qquad C_8=c^2,\qquad L_8=l+bd,\\
 M_8&=j+ax+be+ab^2,\\
 J_1&=\{a,d\},\qquad
 J_2=\{a,e\}+ak,\qquad J_3=\{b,c\}+ak.
\end{split}
 \tag{D.2.19e}
\]
All except $C_8$ are primitive. The reduced diagonal of
$C_8$ is $A_4\otimes A_4$. Each primitive has zero image
in $B$, which has no degree-eight primitive homology.
For the remaining class, $q_*C_8=0$ and
$(q_*\otimes1)\Delta C_8=1\otimes C_8$ show that it is
in the fiber factor. Thus these are actual fiber classes.

Here is the complete operation table needed below:
\[
\begin{array}{c|c|c|c}
 &\beta=\operatorname{Sq}_*^1&\operatorname{Sq}_*^2
                         &\operatorname{Sq}_*^4\\ \hline
 W_8&0&0&0\\
 C_8&0&0&A_4\\
 L_8&Q&d_6&0\\
 M_8&Q&f_6&0\\
 J_1&0&d_6+e_6&0\\
 J_2&0&0&0\\
 J_3&R&d_6+e_6&0
\end{array}
 \tag{D.2.19f}
\]
We give checks fixing the less immediate coefficients.
Serre's cohomology calculation and Cartan give
\[
 \operatorname{Sq}^1(uw)=vw+uv^2,
 \qquad \operatorname{Sq}^2(u^3)=u^4+uv^2.
\]
Consequently
\[
\begin{split}
 \beta L_8&=h+ad=Q,\\
 \beta M_8&=h+a(d+e)+ae=Q,\\
 \operatorname{Sq}_*^2L_8&=b^2=d_6,\\
 \operatorname{Sq}_*^2M_8&=f+ac+a^3=f_6,\\
 \operatorname{Sq}_*^2\{a,d\}
 &=\operatorname{Sq}_*^2\{b,c\}=k=d_6+e_6.
\end{split}
\]
Cartan gives the other first- and second-square entries.
For the last column, instability identifies the fourth
square of the degree-four cohomology class with its cup
square. Primitives annihilate that product, whereas the
diagonal of $C_8$ evaluates to one. This checks the last
column independently of label-operation formulas.
The primary conventions are those of
\cite[\S2, nos.~8--9]{I_Serre1953} and
\cite[Chapter~III, Theorems~3.1--3.2]{I_CohenLadaMay1976}.

The universal coefficient theorem, $H_6F=(\ZZ/2)^2$,
and the two independent first-Bockstein images $Q,R$
in \textup{(D.2.19f)} give
\[
 H_7F=(\ZZ/2)^2.
 \tag{D.2.19g}
\]

\medskip\noindent
\textbf{The entire stable seven-skeleton group.}
Take the finite stable approximation used in D.2.18 through
dimension nine before cancelling odd-unit cell pairs. This
also removes extraneous eight-cells. Write $L_7$ for its
seven-skeleton. Its six-skeleton is
\[
 L_6=C\vee M(\ZZ/4,5)\vee S^6_{\rho_0},
 \qquad C=\cofib\bigl((2,\eta_4):S^4\vee S^5\to S^4\bigr),
\]
and its seven-cell attachments are
\[
 V:\gamma+t,\qquad
 P:\gamma+2\rho_0+\delta,\qquad Q:\delta,\qquad R:0.
 \tag{D.2.19h}
\]
The first three stable stems, including
$\pi_3^S{}_{(2)}=\ZZ/8\{\nu\}$, are recorded in
\cite[\S2, (2.1), p.~149]{I_Oka1986}.
The same Chang cofiber sequence now gives
\[
 0\longrightarrow\ZZ/2\{v_0=i\nu\}
 \longrightarrow\pi_7C
 \longrightarrow\ZZ/2\{\eta_5^2\}\longrightarrow0.
 \tag{D.2.19i}
\]
The class $\delta\eta_6$ maps nontrivially to the last
group and has order two, so this extension splits.
The class $\gamma\eta_6$ maps to zero and is therefore
zero or $v_0$. The fourth-square coefficient from the
bottom four-class into the cone of $v_0$ equals one,
by comparison with the Hopf-invariant-one map $\nu$.
It vanishes for $\gamma\eta_6$: this map factors through
$S^6$, and the induced map from the cone of $\eta_6$
kills the degree-four class while preserving the new
top degree-eight class. Hence
\[
 \gamma\eta_6=0.
 \tag{D.2.19j}
\]

For the Moore summand one has
\[
 \pi_7M(\ZZ/4,5)
 =\ZZ/2\{t\eta_6\}\oplus\ZZ/2\{b_0\},
 \qquad p_*b_0=\eta_6.
 \tag{D.2.19k}
\]
Here the order-two lift can be constructed from the
cofiber map $M(\ZZ/2,5)\to M(\ZZ/4,5)$ that has
degree two on bottom cells and degree one on top cells.
The relation $2\widetilde\eta=i\eta^2$ in the mod-two
Moore spectrum, given in
\cite[\S2, (2.2)(a)--(b), p.~149]{I_Oka1986}, maps to
$2i\eta^2=0$. The top pinch proves independence from
the bottom $\eta^2$ class.
Set $c_0=\rho_0\eta_6$ in the last sphere summand.
The preceding maps induced by \textup{(D.2.19h)} kill
exactly $t\eta_6$ and $\delta\eta_6$ in $\pi_7L_6$.
The kernel of the corresponding map on $\pi_6$ is free
on twice the $Q$ cell and the $R$ cell. The cofiber
sequence therefore gives
\[
 \pi_7L_7=
 \ZZ_{(2)}\{\kappa\}\oplus\ZZ_{(2)}\{\lambda\}
 \oplus(\ZZ/2)\{v_0,b_0,c_0\},
 \quad h(\kappa)=2Q,\quad h(\lambda)=R.
 \tag{D.2.19l}
\]
The extension by the free group splits; the choices of
its lifts are normalized next.

\medskip\noindent
\textbf{Complete detectors for the eight-cell attachments.}
For a map $S^7\to L_7$ with zero mod-two Hurewicz image,
its cone has a new top degree-eight cohomology class.
The second-square coefficients from degree six and the
fourth-square coefficient from degree four are additive
invariants of such stable maps. This follows from
naturality of cones under the stable pinch map.
On the torsion classes they are
\[
\begin{array}{c|c|c}
 &\operatorname{Sq}_*^2&\operatorname{Sq}_*^4\\\hline
 v_0&0&A_4\\
 b_0&d_6+e_6&0\\
 c_0&e_6+f_6&0.
\end{array}
 \tag{D.2.19m}
\]
For $b_0$, collapse the Moore bottom: its top
$\eta$-lift is detected by the six-class $d_6-e_6$.
For $c_0$, use the sphere class $f_6-e_6$. These
factorizations also make their fourth-square coefficients
zero. The coefficient for $v_0$ is the Hopf-invariant-one
coefficient; its second-square coefficients vanish by
the bottom-sphere comparison.

For $\kappa$, the top Bockstein in its cone is $Q$.
The Adem identity
$\operatorname{Sq}^2\operatorname{Sq}^2
=\operatorname{Sq}^3\operatorname{Sq}^1$, together with
$\operatorname{Sq}_*^2Q=z_5$ and $\beta z_5=A_4$,
makes the sum of its three second-square coordinates odd.
Adding $b_0,c_0$ makes this vector exactly $d_6$;
adding $v_0$ makes its fourth-square coefficient zero.
Fix this choice of $\kappa$.
The element $2\lambda$ has zero second- and
fourth-square coefficients: compare its cone with the
mod-two Moore cone of the degree-two map of $S^7$.
The comparison kills the relevant degree-six and
degree-four classes and preserves the top class.

Every eight-cell attachment has even integral degree
on $R$. Its coefficients on $\kappa$ and $\lambda$,
together with \textup{(D.2.19m)}, consequently determine
it completely. The only surviving bottom $\nu$ ambiguity
is the order-two class $v_0$, and its fourth square
detects it. No even multiple of $\nu$ remains hidden.

The integral seven-cycles are free on $Q,R$. By
\textup{(D.2.19g)}, the eight-cell boundary subgroup
is exactly $2\ZZ_{(2)}\{Q,R\}$. As in D.2.18,
choose preimages of $2Q,2R$ congruent to $L_8,J_3$
modulo two, and correct all remaining errors by even
cell combinations. The integral boundaries can thereby
be set to $2Q,2Q,2R$ for $L_8,M_8,J_3$, respectively,
and zero for the other four classes. This preserves
every coefficient in \textup{(D.2.19f)}.
The complete attaching table is
\[
\begin{array}{c|c}
 W_8&0\\
 C_8&v_0\\
 L_8&\kappa\\
 M_8&\kappa+b_0+c_0\\
 J_1&b_0\\
 J_2&0\\
 J_3&2\lambda+b_0.
\end{array}
 \tag{D.2.19n}
\]
Its cokernel is $\ZZ/2\{\lambda\}$, with Hurewicz
image $R$, and it kills $b_0$ and $c_0$ separately.
Cells of dimension nine or more do not change this
stable group. Equation \textup{(D.2.19c)} proves
\textup{(D.2.19a)} for the actual unstable group.

The inclusion of the Moore summand sends its top
$\eta$-lift to $b_0=0$. Any other Moore extension with
the same bottom differs by an element of $\pi_6F$
followed by the top pinch. Composing that difference
with a top $\eta$-lift gives a multiple of
$\rho\eta_6=0$. The actual extension $\mu$ of D.2.17
therefore satisfies \textup{(D.2.19b)}, independently
of its permitted top-cell choice. These statements
also hold in $T$, since $\pi_iF\cong\pi_iT$ in the
degrees used here.

\medskip\noindent
\textbf{The entire source group in the next attaching dimension.}
Write $Q_i=K_i/K_4$, with the source skeleta of D.2.18.
After the indicated cell cancellations, its seven-dimensional
quotient has the following two-local model:
\[
 Q_7\simeq D\vee M(\ZZ/2,6)_z\vee S^7_F\vee S^7_J,
 \qquad
 D=\operatorname{cofib}\left(
 \begin{pmatrix}4&\eta_5\\0&2\end{pmatrix}:
 S^5\vee S^6\longrightarrow S^5\vee S^6\right).
 \tag{D.2.19o}
\]
The six-dimensional coordinates are the Moore top $d=D$,
the free cycle $e=D+E$, and $z=A_0^2$ before their seven-cell
attachments. The free seven-classes are
$F=\sigma f$ and $J=A_0B_0+Q^4A_0$ modulo two.
We distinguish the complex $D$ in \textup{(D.2.19o)} from
the homology class $D=\sigma d$ by its argument whenever
necessary.

There is a complete and explicitly detected group
\[
 \pi_7Q_7=
 \ZZ_{(2)}\{F,J\}\oplus(\ZZ/2)\{\tau_d,\tau_e,\tau_z\}.
 \tag{D.2.19p}
\]
To see this without presuming a Moore-group extension,
collapse the bottom five-sphere of the complex $D$.
Its quotient is $S^6_d\vee M(\ZZ/2,6)_e$.
The map $\pi_7S^5\to\pi_7D$ is zero: its generator
$i\eta_5\eta_6$ is
$(i\eta_5+2e)\eta_6$, since $2\eta_6=0$.
Thus the quotient injects $\pi_7D$ into
$\pi_7(S^6\vee M(\ZZ/2,6))=(\ZZ/2)^2$.
The connecting map to $\pi_6S^5$ vanishes on both generators.
On $S^6_d$ it is the degree-four map, which kills $\eta$;
on the Moore factor it is $\eta_6$ followed after its top
pinch, which kills its bottom $\eta$ generator. Therefore
the injection is an isomorphism.

The image of any lift
$\zeta:S^7\to M(\ZZ/4,5)$ with $p\zeta=\eta_6$
is the class $\tau_d$. Its choice ambiguity is $i\eta^2$,
which has just been killed in $D$.
The other two classes are
$\tau_e=e\eta_6$ and $\tau_z=z\eta_6$.
They are detected exactly by the second squares from the
three degree-six cohomology classes $d^\vee,e^\vee,z^\vee$.
The remaining two components in \textup{(D.2.19p)} are
their integral degrees on $F,J$.
The quotient is four-connected, so these stable groups and
attaching components are the actual unstable ones. In
particular, there is no additional Whitehead component.

\medskip\noindent
\textbf{Every degree-eight source class, including the products.}
The full ordinary mod-two basis is
\[
 G_8=\sigma g,\quad H_8=\sigma h,\quad
 A_0C,\quad B_0^2,\quad Q^5A_0.
 \tag{D.2.19q}
\]
The two unit classes are dual to the suspended monomials
$u^2v,uw$. The products and the new arity-two class are
included; arities at least three start in dimension nine.
The complete operations relevant to
\textup{(D.2.19p)} are
\[
\begin{array}{c|c|c}
 \text{class}&\beta&\operatorname{Sq}_*^2\\ \hline
 G_8&F&E=d+e\\
 H_8&0&E=d+e\\
 A_0C&0&z\\
 B_0^2&0&z\\
 Q^5A_0&0&z.
\end{array}
 \tag{D.2.19r}
\]
For the unit classes use
$\operatorname{Sq}^2w=0$ and
$\operatorname{Sq}^2(uv)=u^2v+uw$.
Cartan gives
$\operatorname{Sq}_*^2(A_0C)=A_0^2$ and
$\operatorname{Sq}_*^2(B_0^2)=(\beta B_0)^2=A_0^2$.
The stable Nishida formula gives
\[
 \operatorname{Sq}_*^2Q^5A_0
   =\binom{3}{2}Q^3A_0=A_0^2,
 \tag{D.2.19s}
\]
because $\beta A_0=0$. These are the primary Serre and
Cohen--Lada--May conventions already cited above.

\medskip\noindent
\textbf{The degree-four boundary coefficient is an integral calculation.}
The Snaith splitting identifies
\[
 H_7(A;\ZZ)=\ZZ/2\{F\}\oplus\ZZ/4\{J\}.
 \tag{D.2.19t}
\]
We verify the exponent and the particular generator of its
second summand. For the first summand, Serre's mod-two
dimensions and the universal coefficient theorem make
$H_6(B;\ZZ)$ cyclic two-primary. Its first Bockstein
$\beta g=f\ne0$ fixes its exponent as two.

For weight two, use the bottom Moore chains of $X$:
$a$ has degree three, $b$ has degree four, and $db=2a$.
Use free $E\Sigma_2$ cells with
$de_1=(\tau-1)e_0$ and $de_2=(1+\tau)e_1$.
The Koszul action satisfies
$\tau(a\otimes a)=-a\otimes a$ and
$\tau(a\otimes b)=b\otimes a$.
Set
\[
\begin{gathered}
 x_0=e_0(a\otimes a),\quad r_1=e_1(a\otimes a),\quad
 t=e_0(a\otimes b),\quad t'=e_0(b\otimes a),\\
 h=e_1(a\otimes b),\qquad q_0=e_0(b\otimes b).
\end{gathered}
\]
The degree-seven boundaries are
\[
 dr_1=-2x_0,\qquad dt=-2x_0,\qquad dt'=2x_0.
\]
Hence $w_0=t-r_1$ and $w_0'=t'+r_1$ freely generate
the integral seven-cycles. The next boundaries are
\[
 dh=w_0'-w_0,\qquad dq_0=2(w_0+w_0'),
 \qquad d(e_2(a\otimes a))=0.
 \tag{D.2.19u}
\]
The other mixed $e_1$ cell has the opposite boundary.
Higher-label cells in dimension eight are cycles and do
not alter the degree-seven group. Thus its weight-two part is
\[
 \frac{\ZZ\{w_0,w_0'\}}
      {\langle w_0'-w_0,\ 2(w_0+w_0')\rangle}
 =\ZZ/4\{w_0\}.
\]
Modulo two, $w_0$ is exactly $J=A_0B_0+Q^4A_0$.
Furthermore
\[
 d(q_0-2h)=4w_0,
 \qquad \beta_2[B_0^2]=[J]\ne0.
 \tag{D.2.19v}
\]
This uses a corrected integral representative; $dq_0$
itself is not divisible by four.
The cells $A_0C$ and $Q^5A_0$ have integral cycle
representatives. The unit class $H_8$ has no
second-Bockstein component on the weight-two class $J$.
The arity separation is provided by
\cite[Theorem~1.1 and \S2]{Snaith1974}; the displayed
chain formulas use the conventions of
\cite[Chapter~III, proof of Lemma~3.10, p.~232]{I_CohenLadaMay1976}.
They prove \textup{(D.2.19t)} and the coefficient in
\textup{(D.2.19v)}, rather than inferring that coefficient
from a mod-two dimension count.

\medskip\noindent
\textbf{The actual eight-cell matrix.}
The free seven-cycles of $Q_7$ are $F,J$. Equation
\textup{(D.2.19t)} makes their boundary sublattice
$\langle2F,4J\rangle$. The first and second Bocksteins
specify a normal form preserving all the displayed
mod-two cell labels. In detail, a term $4mJ$ in the
boundary of $G_8$ is removed by replacing the free
generator $F$ by $F+2mJ$, with unchanged mod-two class.
Terms $4mF$ in the boundary of $B_0^2$ are removed by
even multiples of $G_8$. On a zero-first-Bockstein
cell the remaining $F$ coefficient is even, and a
zero second Bockstein on $J$ makes the remaining
$J$ coefficient an even multiple of the degree-four
relation. Even cell combinations therefore remove
these terms without changing a square coefficient.
Odd rescalings are invertible two-locally. This gives
boundaries $2F$ and $4J$ for $G_8$ and $B_0^2$,
and zero integral degrees for the other three cells.

The detectors in \textup{(D.2.19p)} and the complete
table \textup{(D.2.19r)} now give every attachment:
\[
\begin{array}{c|c}
 \text{eight-cell}&\text{attaching map in }\pi_7Q_7\\ \hline
 G_8&2F+\tau_d+\tau_e\\
 H_8&\tau_d+\tau_e\\
 B_0^2&4J+\tau_z\\
 A_0C&\tau_z\\
 Q^5A_0&\tau_z.
\end{array}
 \tag{D.2.19w}
\]
The calculation may be made after cancelling odd unit
pairs in a two-local finite model. Each comparison is
realizable unstably here: the quotients are four-connected,
their relevant attaching dimensions are at most seven,
and Freudenthal identifies these stable attaching maps
with ordinary ones. Pullback along localization gives
ordinary maps on the original finite source skeleta,
because the target $T$ is two-local. Retaining the
cancelled pairs gives the same mapping problem.

The correction of D.2.18 is $\mu$ on the Moore
factor, $\rho$ on the six-sphere $e$, and zero on
$z$. Choose its values on the free seven-spheres
$F,J$ to be zero. Its values on the torsion classes are
\[
 \tau_d\longmapsto\mu\zeta=0,\qquad
 \tau_e\longmapsto\rho\eta_6=0,\qquad
 \tau_z\longmapsto0.
 \tag{D.2.19x}
\]
The first two vanish by the target calculation.
They are independent of the chosen nullhomotopy on
the seven-cell $i\eta_5+2e$, because the two
representing maps already lie in its six-skeleton.
The ambiguity $i\eta^2$ in $\zeta$ is also harmless:
$\overline\alpha\eta^2=2\rho\eta_6=0$.
Thus every attachment in \textup{(D.2.19w)} is killed,
and the correction extends over $Q_8$.

\medskip\noindent
\textbf{The canonical summand also extends to the new source cells.}
The canonical map $g_2:C_2X\to T$ must not be treated
as already defined on the new eight-cell of $A$.
We verify the needed extension independently.
Choose a bottom generator $i:S^3\to X$ and put
$X'=X/S^3$. This quotient is three-connected:
$X$ is two-connected, and
$H_3(S^3;\ZZ)\to H_3X=\ZZ/2$ is surjective.
The Hurewicz and cofiber sequences give
$H_3X'=0$, hence $\pi_3X'=0$.

Freudenthal makes $B\to\Omega^2\Sigma^2B$
three-connected. Since $\pi_3B=0$, it follows that
$\pi_5Y=0$, as already checked in D.2.13.
Consequently
$h_0\Sigma^2i:S^5\to Y$ is null; equivalently
its bottom description is $u_4\eta_4$, in this
zero homotopy group. The cofiber universal property
therefore gives a pointed map
\[
 \overline h:\Sigma^2X'\longrightarrow Y,
 \qquad \overline h\Sigma^2q\simeq h_0,
 \tag{D.2.19y}
\]
where $q:X\to X'$ is the quotient.

Since $X'$ is three-connected, the stabilization
$C_2X'\to QX'$ is nine-connected: the unit on
the five-connected space $\Sigma^2X'$ is
eleven-connected, two loops subtract two, and
subsequent stages improve the bound.
Form the unsuspended homotopy pullback
\[
 P_0=(QX)\times^h_{QX'} C_2X'
          \longrightarrow QX.
 \tag{D.2.19z}
\]
It is nine-connected and receives the natural lift
of $C_2X$. Projection to $C_2X'$ followed by
$\Omega^2\overline h$ gives $P_0\to T$; its
restriction to $C_2X$ is $g_2$, by
\textup{(D.2.19y)} and naturality.
After fibration replacement the fiber over $QX$
is eight-connected. The specified lift therefore
extends across every relative cell of dimension
at most nine. This gives a canonical summand on
the required $K_8$ (indeed through source relative
dimension nine), with its original restriction
and low homology values.

Finally pull the extended correction on $Q_8$
back along $K_8\to Q_8$ and add it to this
canonical summand using multiplication in $T$.
The result is an actual ordinary map $G_8:K_8\to T$.
Its restriction through dimension six is the
chosen map of D.2.17. The argument treats the new
eight-cell explicitly and uses the unsuspended
pullback \textup{(D.2.19z)}; it does not assume
that suspension preserves a homotopy fiber sequence.
\end{proof}
